\section{Proofs of Theoretical Results}
\label{sec:proofs}

\subsection{Proof of \texorpdfstring{\cref{prop:offline}}{Proposition (Offline)}}
\label{sec:proof-offline}

We prove the three parts in turn.

\paragraph{Part~(i): prompt-independence.}
Under $\mathcal{D}_t(c,x_t)=\mathcal{D}_t$,
$G_t(K_t)$ depends only on $(t,K_t)$, so
$\sum_t G_t(K_t)$ is the same function of $\{K_t\}$ for every $(c,x)$.
Any $\{K_t^*\}$ solving~\eqref{eq:alloc} is therefore simultaneously
optimal for every prompt.

\paragraph{Part~(ii): water-filling monotonicity.}
First, treat $K$ as continuous.
Under the location-scale assumption~\eqref{eq:gain-sde},
$G_t(K)=\sigma_t\,a_K$ is strictly concave in $K$
(since $a_K$ is strictly concave; see David and Nagaraja, 2003,
Ch.~4 for the Gaussian case; the integer solution is obtained by
rounding and loses at most one unit of gain per timestep).
Then, write the Lagrangian of~\eqref{eq:alloc} as
$\mathcal{L}=\sum_t G_t(K_t)-\lambda\bigl(\sum_t K_t - B\bigr)
+\sum_t \mu_t(K_t-1)$,
which yields the first-order and complementary-slackness conditions
\[
  G_t'(K_t^*)=\lambda^*-\mu_t^*,
  \qquad
  \mu_t^*\ge 0,
  \qquad
  \mu_t^*(K_t^*-1)=0.
\]
If $K_t^*>1$ then $\mu_t^*=0$ and the discrete marginal gain equals the
shadow price: $\sigma_t\,\Delta a_{K_t^*}=\lambda^*$.
Hence
\[
  \Delta a_{K_t^*} = \frac{\lambda^*}{\sigma_t}.
\]
Because the location-scale assumption makes $a_K$ the \emph{same}
function for every~$t$, the right-hand side pins $K_t^*$ on a single
strictly decreasing sequence $\Delta a_K$.
Hence $\sigma_{t_1}>\sigma_{t_2}$
implies $\Delta a_{K_{t_1}^*} < \Delta a_{K_{t_2}^*}$, hence
$K_{t_1}^*>K_{t_2}^*$.
(If the shape of $\mathcal{D}_t$ varied with~$t$, each timestep would
have its own curve $a_{K,t}'$ and this monotonicity would not follow
from $\sigma_{t_1}>\sigma_{t_2}$ alone.)

\paragraph{Part~(iii): strict suboptimality of uniform allocation.}
At the uniform point $K_t\equiv B/T$, the marginal gains are
$G_t'(B/T)=\sigma_t\,a_{B/T}'$.
If $\sigma_{t_1}\neq\sigma_{t_2}$ for some pair, these marginals differ,
violating the KKT condition.
To confirm a strict improvement exists, consider the budget-neutral
perturbation $K_{t_1}=B/T+\epsilon$,
$K_{t_2}=B/T-\epsilon$ with all other $K_t$ unchanged.
Expanding to first order,
\begin{align*}
  &\;\bigl[G_{t_1}(B/T+\epsilon)+G_{t_2}(B/T-\epsilon)\bigr]
   -\bigl[G_{t_1}(B/T)+G_{t_2}(B/T)\bigr] \\
  &= \epsilon\bigl[G_{t_1}'(B/T)-G_{t_2}'(B/T)\bigr]+O(\epsilon^2) \\
  &= \epsilon\,(\sigma_{t_1}-\sigma_{t_2})\,a_{B/T}'+O(\epsilon^2)
   \;>\;0
\end{align*}
for small $\epsilon>0$.
The common factor $a_{B/T}'$ can be extracted because $a_K$ is
$t$-independent (location-scale); combined with
$\sigma_{t_1}-\sigma_{t_2}>0$ and $a_{B/T}'>0$, the perturbation
is strictly positive.
The uniform allocation is therefore not a stationary point.
Since the objective is concave over the convex constraint set
$\{\sum_t K_t=B,\;K_t\ge 1\}$, any non-stationary feasible point
is strictly suboptimal, so the global maximizer $\{K_t^*\}$
must satisfy $\sum_t G_t(K_t^*)>\sum_t G_t(B/T)$ whenever
$\{\sigma_t\}$ are not all equal.
\hfill$\blacksquare$

\subsection{Proof of \texorpdfstring{\cref{prop:online}}{Proposition (Online)}}
\label{sec:proof-online}

\paragraph{Part~(i): convexity and Jensen gap.}
For fixed $\{K_t\}$ with $\sum_t K_t=B$ and $K_t\ge 1$, the map
$\boldsymbol{\sigma}\mapsto\sum_t\sigma_t\,a_{K_t}$ is linear.
Denote the \emph{instance-optimal value}, the best total gain attainable
when the true sensitivity vector $\boldsymbol{\sigma}$ is known, by
\[
  V^*(\boldsymbol{\sigma})
  \;=\;
  \max_{\{K_t\}:\,\sum_t K_t=B,\;K_t\ge 1}
  \sum_t\sigma_t\,a_{K_t}
  \;=\;
  \sup_{\{K_t\}\in\mathcal{K}}\;
  \langle\boldsymbol{a}_K,\,\boldsymbol{\sigma}\rangle.
\]
This is a pointwise supremum of linear functions, therefore convex
(Boyd and Vandenberghe, 2004, \S3.2.3).
By Jensen's inequality,
\[
  \mathbb{E}_{\boldsymbol{\sigma}}
  \bigl[V^*(\boldsymbol{\sigma})\bigr]
  \;\ge\;
  V^*\!\bigl(\mathbb{E}[\boldsymbol{\sigma}]\bigr)
  \;=\;
  V^*(\bar{\boldsymbol{\sigma}}).
\]
Since $\mathcal{K}$ is finite (integer allocations), $V^*$ is the
maximum of finitely many linear functions, hence convex and
piecewise-linear.
Equality holds whenever there exists an allocation
$K^\dagger\in\mathcal{K}$ such that
$V^*(\boldsymbol{\sigma})=\sum_t \sigma_t a_{K_t^\dagger}$ almost
surely on the support of $\boldsymbol{\sigma}$.
In particular, a constant variance profile is a sufficient special case.

\paragraph{Part~(ii): dual-threshold stopping.}
Under the location-scale model~\eqref{eq:gain-sde},
\[
  \Delta_t(K)
  := G_t(K+1)-G_t(K)
  = \sigma_t\,(a_{K+1}-a_K) + O_{K,d,L_t}(g_t^2 h)
\]
gives the \emph{discrete marginal gain}: the additional value from one
more search iteration at timestep~$t$.
The leading term separates into
\[
  \sigma_t\cdot \Delta a_K,
  \qquad
  \Delta a_K:=a_{K+1}-a_K .
\]
Because $a_K$ is $t$-independent, this marginal gain separates into a
timestep-specific sensitivity~$\sigma_t$ and a universal exhaustion
sequence~$\Delta a_K$; this separation is what allows the two
thresholds below to monitor each factor independently.

Next, characterize the optimal stopping rule.
Search at timestep $t$ should stop when $\Delta_t(K)\le\lambda^*$, where
$\lambda^*$ is the shadow price of the budget constraint.
For each completed step $s$, define the \emph{average realized gain}
\[
  \bar{g}_s
  := \frac{G_s(\hat{K}_s)-G_s(1)}{\hat{K}_s - 1}
  = \frac{1}{\hat{K}_s-1}
    \sum_{k=1}^{\hat{K}_s-1}\Delta_s(k).
\]
Concavity of $G_s$ implies that the discrete marginals
$\Delta_s(k)$ are non-increasing in~$k$, so the average dominates the
last active marginal:
\[
  \bar{g}_s
  \ge \Delta_s(\hat{K}_s-1).
\]
Averaging over the gain-history window $\mathcal{H}_g$:
\[
  \bar{g}
  =\frac{1}{|\mathcal{H}_g|}\sum_s\bar{g}_s
  \;\ge\;
  \frac{1}{|\mathcal{H}_g|}\sum_s \Delta_s(\hat{K}_s-1).
\]
If the realized allocations $\{\hat{K}_s\}$ coincide with the optimal
solution of~\eqref{eq:alloc}, then
$\Delta_s(\hat{K}_s-1)\ge\lambda^*\ge \Delta_s(\hat{K}_s)$ for every
active $s$ by the discrete KKT conditions for separable concave
allocation, giving $\bar{g}\ge\lambda^*$.
In practice $\{\hat{K}_s\}$ are determined online and need not be
exactly optimal, so $\beta_g\bar{g}$ ($\beta_g\in(0,1)$) serves as a
conservative, adaptive estimate of $\lambda^*$.

Finally, explain why both thresholds are needed.
The gain threshold $\widetilde{g}_t^{(j)}<\beta_g\bar{g}$ tests the
leading-order product $\sigma_t\,\Delta a_K$ against the
shadow-price scale~$\lambda^*$.
By \cref{thm:loc-scale}(ii), the conditional score variance obeys
$\operatorname{Var}(S_t\mid X_t,c)=\sigma_t^2+O_d(g_t^3 h^{3/2})$, so a
low empirical within-timestep score variance is a natural proxy for low
sensitivity, but not an exact test of $\sigma_t$.
Requiring both conditions before stopping avoids
premature termination at high-value timesteps under estimation noise.
If only the gain threshold is checked, a timestep with high
$\sigma_t$ but temporarily low $\Delta a_K$ (exhaustion from a few
unlucky draws) would be abandoned even though a single additional
candidate could yield a large jump.
Conversely, if only the variance threshold proxy is checked, a low-$\sigma_t$
timestep would be abandoned even when its $\Delta a_K$ is still large and
the search has not yet diminished.
The conjunction ensures that stopping occurs only when both the
timestep's intrinsic sensitivity and its remaining search potential
are genuinely low, preventing the loss of high-value search
opportunities.
\hfill$\blacksquare$
